\documentclass[12pt]{article}
\usepackage[margin=1in, top=0.75in, bottom=0.75in]{geometry}
% worksheet-preamble.tex — the shipped preamble every generated document
% \input's. SOURCE OF TRUTH: this file. references/latex-templates.md documents
% these macros but no longer serves as the text agents transcribe — retyping a
% ~130-line preamble per document, three documents per run, is exactly the
% transcription-drift error class the rest of the pipeline exists to catch.
%
% Usage (compile.sh stages this file beside your .tex automatically):
%   \documentclass[12pt]{article}
%   \usepackage[margin=1in, top=0.75in, bottom=0.75in]{geometry}
%   % worksheet-preamble.tex — the shipped preamble every generated document
% \input's. SOURCE OF TRUTH: this file. references/latex-templates.md documents
% these macros but no longer serves as the text agents transcribe — retyping a
% ~130-line preamble per document, three documents per run, is exactly the
% transcription-drift error class the rest of the pipeline exists to catch.
%
% Usage (compile.sh stages this file beside your .tex automatically):
%   \documentclass[12pt]{article}
%   \usepackage[margin=1in, top=0.75in, bottom=0.75in]{geometry}
%   % worksheet-preamble.tex — the shipped preamble every generated document
% \input's. SOURCE OF TRUTH: this file. references/latex-templates.md documents
% these macros but no longer serves as the text agents transcribe — retyping a
% ~130-line preamble per document, three documents per run, is exactly the
% transcription-drift error class the rest of the pipeline exists to catch.
%
% Usage (compile.sh stages this file beside your .tex automatically):
%   \documentclass[12pt]{article}
%   \usepackage[margin=1in, top=0.75in, bottom=0.75in]{geometry}
%   \input{worksheet-preamble}
%   \input{figure-macros}          % when using the shipped figure macros
%
% geometry stays in the document, NOT here: worksheets use margin=1in
% (top/bottom 0.75in) while study guides use margin=0.85in (top/bottom 0.7in),
% and geometry options cannot be re-issued once the package is loaded.
%
% ASCII only in this file: pdflatex (the fallback engine) cannot typeset
% literal Unicode symbols, so documents must compile identically under both
% tectonic and pdflatex.

% ── Required packages (superset of the worksheet and study-guide shells; a
% worksheet loading mdframed it does not use costs nothing, one shared file
% that cannot drift is worth it) ─────────────────────────────────────────────
\usepackage{amsmath, amssymb}
\usepackage{tikz, pgfplots}
\usepackage{enumitem, fancyhdr, multicol, array, booktabs}
\usepackage{mdframed, xcolor}
\pgfplotsset{compat=1.18}
\usetikzlibrary{calc, angles, quotes}   % needed by the geometric figure macros

% ── Problem numbering ───────────────────────────────────────────────────────
\newcounter{prob}
\newcommand{\problem}[1]{\stepcounter{prob}\vspace{0.4cm}\noindent\textbf{\large\theprob.}\quad #1\vspace{0.3cm}}

% ── Shrink-to-fit title: \LARGE when it fits on one line, otherwise scale
% down to the text width — never enlarges, never wraps mid-phrase. (\resizebox
% comes from graphicx, which tikz already loads.) ────────────────────────────
\newlength{\fittedw}
\newcommand{\fittedtitle}[1]{%
  \settowidth{\fittedw}{\LARGE\bfseries #1}%
  \ifdim\fittedw>\linewidth\resizebox{\linewidth}{!}{\bfseries #1}%
  \else{\LARGE\bfseries #1}\fi}

% ── Running headers. Keep the left title SHORT (<= ~28 chars, e.g. "Triangle
% Trig Practice", not the full course name) — it shares the header line with
% the Name/Date blanks and a long title overlaps them. ──────────────────────
% \wsheader{Short Title} — student worksheet
\newcommand{\wsheader}[1]{%
  \pagestyle{fancy}\fancyhf{}%
  \fancyhead[L]{\textbf{\small #1}}%
  \fancyhead[R]{\small Name: \underline{\hspace{4.5cm}}~Date: \underline{\hspace{1.8cm}}}%
  \fancyfoot[C]{\thepage}%
  \renewcommand{\headrulewidth}{0.4pt}}

% \akheader{Short Topic} — answer key
\newcommand{\akheader}[1]{%
  \pagestyle{fancy}\fancyhf{}%
  \fancyhead[L]{\textbf{\small #1 --- Answer Key}}%
  \fancyhead[R]{\textit{\small For instructor/parent use}}%
  \fancyfoot[C]{\thepage}%
  \renewcommand{\headrulewidth}{0.4pt}}

% \ssheader{Topic} — skills summary / study guide
\newcommand{\ssheader}[1]{%
  \pagestyle{fancy}\fancyhf{}%
  \fancyhead[L]{\textbf{Skills Summary: #1}}%
  \fancyhead[R]{\small\textit{Study Guide \& Reference}}%
  \fancyfoot[C]{\thepage}%
  \renewcommand{\headrulewidth}{0.4pt}}

% ── Title blocks ────────────────────────────────────────────────────────────
% \wstitleblock{Title}{Course}{Date}
\newcommand{\wstitleblock}[3]{%
  \begin{center}
    \fittedtitle{#1}\\[0.3cm]
    {\large #2 \quad $\bullet$ \quad #3}
  \end{center}
  \noindent\rule{\linewidth}{0.4pt}\vspace{0.2cm}}

% \aktitleblock{Topic}{Course}{Date} — "Answer Key" is a subtitle BENEATH the
% topic, never appended to it: a long topic would wrap mid-phrase.
\newcommand{\aktitleblock}[3]{%
  \begin{center}
    \fittedtitle{#1}\\[0.15cm]
    {\Large\bfseries Answer Key}\\[0.25cm]
    {\large #2 \quad $\bullet$ \quad #3}
  \end{center}}

% \sstitleblock{Topic}
\newcommand{\sstitleblock}[1]{%
  \begin{center}
    {\LARGE\textbf{Skills Summary}}\\[0.2cm]
    {\large\textbf{#1}}\\[0.1cm]
    {\small\textit{Use this reference while completing your worksheet or when studying.}}
  \end{center}
  \vspace{0.1cm}
  \noindent\rule{\linewidth}{1.5pt}
  \vspace{0.3cm}}

% ── Study-guide color palette and boxes ─────────────────────────────────────
\definecolor{skillblue}{RGB}{30,100,180}
\definecolor{skillbluebg}{RGB}{235,244,255}
\definecolor{warnorange}{RGB}{200,90,0}
\definecolor{warnbg}{RGB}{255,243,230}
\definecolor{exgreen}{RGB}{20,120,60}
\definecolor{exgreenbg}{RGB}{230,248,238}

% Formula/rule box
\newmdenv[
  backgroundcolor=skillbluebg,
  linecolor=skillblue, linewidth=1.5pt,
  innertopmargin=6pt, innerbottommargin=6pt,
  innerleftmargin=10pt, innerrightmargin=10pt,
  skipabove=6pt, skipbelow=4pt
]{formulabox}

% Mini example box
\newmdenv[
  backgroundcolor=exgreenbg,
  linecolor=exgreen, linewidth=1pt,
  innertopmargin=5pt, innerbottommargin=5pt,
  innerleftmargin=10pt, innerrightmargin=10pt,
  skipabove=4pt, skipbelow=4pt
]{examplebox}

% Watch-out box
\newmdenv[
  backgroundcolor=warnbg,
  linecolor=warnorange, linewidth=1pt,
  innertopmargin=5pt, innerbottommargin=5pt,
  innerleftmargin=10pt, innerrightmargin=10pt,
  skipabove=4pt, skipbelow=4pt
]{watchoutbox}

% Skill section heading
\newcommand{\skillheading}[1]{%
  \vspace{0.4cm}
  {\large\textbf{\textcolor{skillblue}{#1}}}
  \vspace{0.1cm}
  \hrule height 1pt
  \vspace{0.2cm}
}

%   % figure-macros.tex — shipped geometric figure macros. SOURCE OF TRUTH: this
% file; references/latex-templates.md documents usage. Requires the tikz
% libraries loaded by worksheet-preamble.tex (calc, angles, quotes), so
% \input worksheet-preamble first.
%
% CONTRACT with the layout/prose checkers (do not break it):
%   - Macro names contain rt/tri/fig — that is how tests/check_layout.py and
%     tests/check_prose_consistency.py recognize a macro-style figure.
%   - Mandatory arguments are FLAT braces only (no nested {} inside an arg);
%     the checkers' arg scanners match [^{}]* per group.
%   - Value-bearing args carry exactly the printed label text; scale/styling
%     goes in the single optional [..] argument, which the checkers ignore
%     (so scale=1.2 is never mistaken for a figure value).
%
% Figure conventions (from SKILL.md / latex-templates.md): every number
% printed in a figure must come from the problem's verify JSON; vertices are
% A/B/C with sides a/b/c OPPOSITE them — the same convention verify.py's
% triangle type uses. Right triangles put the right angle at C, so c (= AB)
% is the hypotenuse — the SAME convention scripts/render_figures.py constructs
% (its RIGHT_ANGLE_VERTEX constant): \refrt sits beside renderer-built
% \probfig figures on the page, so the two sources must never disagree.
% tests/test_figure_convention.py enforces this by parsing the
% "% right angle mark" lines below — that trailing comment is load-bearing;
% keep it on every right-angle-mark \draw.
%
% ASCII only: pdflatex (the fallback engine) cannot typeset literal Unicode.

% ── \rtfig[scale]{bottom-leg label}{right-leg label}{hypotenuse label} ──────
% Right triangle, right angle at C (hypotenuse c = AB — the renderer's
% convention). A=(0,0), C=(4,0), B=(4,3): the 3-4-5 shape keeps every label
% clear of the drawn lines at default scale. Bottom leg = b (AC), right leg
% = a (CB), hypotenuse = c (AB) — name sides accordingly in the labels.
% Example: \rtfig{$8$}{$6$}{$x$}   or   \rtfig[0.9]{$b = 8$}{$a = 6$}{$c$}
\newcommand{\rtfig}[4][1.2]{%
  \begin{center}
  \begin{tikzpicture}[scale=#1]
    \coordinate (A) at (0,0);
    \coordinate (B) at (4,3);
    \coordinate (C) at (4,0);
    \draw[thick] (A) -- (C) -- (B) -- cycle;
    \draw (C) ++(-.25,0) -- ++(0,.25) -- ++(.25,0);  % right angle mark
    \node[below left] at (A) {$A$};
    \node[above right] at (B) {$B$};
    \node[below right] at (C) {$C$};
    \node[below] at (2,0) {#2};
    \node[right] at (4,1.5) {#3};
    \node[above left] at (2,1.5) {#4};
  \end{tikzpicture}
  \end{center}}

% ── \trifig[scale]{c}{b}{A-deg}{c-label}{b-label}{a-label}{angle-label} ─────
% General triangle, TO SCALE by construction (SAS): A at the origin, B on the
% x-axis at distance c, C computed from side b and angle A (TikZ trig takes
% degrees). Numeric args #2-#4 are the ACTUAL given values, so the drawing
% cannot disagree with the math; label args #5-#8 carry what is printed.
% Example: \trifig{7}{5}{34}{$c = 7$}{$b = 5$}{$a = ?$}{$34^\circ$}
\newcommand{\trifig}[8][0.6]{%
  \begin{center}
  \begin{tikzpicture}[scale=#1]
    \coordinate (A) at (0,0);
    \coordinate (B) at (#2,0);
    \coordinate (C) at ({#3*cos(#4)},{#3*sin(#4)});
    \draw[thick] (A) -- (B) -- (C) -- cycle;
    \node[below left]  at (A) {$A$};
    \node[below right] at (B) {$B$};
    \node[above]       at (C) {$C$};
    \node[below]       at ($(A)!0.5!(B)$) {#5};
    \node[above left]  at ($(A)!0.5!(C)$) {#6};
    \node[above right] at ($(B)!0.5!(C)$) {#7};
    \pic [draw, angle radius=7mm, angle eccentricity=1.6, "#8"] {angle = B--A--C};
  \end{tikzpicture}
  \end{center}}

% ── \refrt — the value-free reference figure ────────────────────────────────
% SKILL.md's figure-scope rule requires shared labelling conventions to live
% in ONE value-free reference figure placed with the directions; until now no
% template for it existed, so every generation invented one — and an invented
% figure that accidentally carried a number would itself trip the
% all-or-nothing figure rule it exists to satisfy. This macro has zero
% numerals in its output by construction (it takes no arguments and prints
% only letters), and the caption is baked in verbatim so the figure cannot be
% mistaken for a problem's givens. Right angle at C, so sides read: a = BC
% (opposite A), b = AC (opposite B), c = AB (opposite C, the hypotenuse) —
% matching render_figures.py's RIGHT_ANGLE_VERTEX, because this reference
% figure defines the convention students apply to every renderer-built
% figure on the sheet.
\newcommand{\refrt}{%
  \begin{center}
  \begin{tikzpicture}[scale=1.0]
    \coordinate (A) at (0,0);
    \coordinate (B) at (3.6,2.7);
    \coordinate (C) at (3.6,0);
    \draw[thick] (A) -- (C) -- (B) -- cycle;
    \draw (C) ++(-.25,0) -- ++(0,.25) -- ++(.25,0);  % right angle mark
    \node[below left] at (A) {$A$};
    \node[above right] at (B) {$B$};
    \node[below right] at (C) {$C$};
    \node[below] at (1.8,0) {$b$};
    \node[right] at (3.6,1.35) {$a$};
    \node[above left] at (1.8,1.35) {$c$};
    \pic [draw, angle radius=6mm, angle eccentricity=1.7, "$A$"] {angle = C--A--B};
  \end{tikzpicture}\\
  {\small\textit{How every triangle here is labelled. No values shown: use the
  numbers given in each problem.}}
  \end{center}}
          % when using the shipped figure macros
%
% geometry stays in the document, NOT here: worksheets use margin=1in
% (top/bottom 0.75in) while study guides use margin=0.85in (top/bottom 0.7in),
% and geometry options cannot be re-issued once the package is loaded.
%
% ASCII only in this file: pdflatex (the fallback engine) cannot typeset
% literal Unicode symbols, so documents must compile identically under both
% tectonic and pdflatex.

% ── Required packages (superset of the worksheet and study-guide shells; a
% worksheet loading mdframed it does not use costs nothing, one shared file
% that cannot drift is worth it) ─────────────────────────────────────────────
\usepackage{amsmath, amssymb}
\usepackage{tikz, pgfplots}
\usepackage{enumitem, fancyhdr, multicol, array, booktabs}
\usepackage{mdframed, xcolor}
\pgfplotsset{compat=1.18}
\usetikzlibrary{calc, angles, quotes}   % needed by the geometric figure macros

% ── Problem numbering ───────────────────────────────────────────────────────
\newcounter{prob}
\newcommand{\problem}[1]{\stepcounter{prob}\vspace{0.4cm}\noindent\textbf{\large\theprob.}\quad #1\vspace{0.3cm}}

% ── Shrink-to-fit title: \LARGE when it fits on one line, otherwise scale
% down to the text width — never enlarges, never wraps mid-phrase. (\resizebox
% comes from graphicx, which tikz already loads.) ────────────────────────────
\newlength{\fittedw}
\newcommand{\fittedtitle}[1]{%
  \settowidth{\fittedw}{\LARGE\bfseries #1}%
  \ifdim\fittedw>\linewidth\resizebox{\linewidth}{!}{\bfseries #1}%
  \else{\LARGE\bfseries #1}\fi}

% ── Running headers. Keep the left title SHORT (<= ~28 chars, e.g. "Triangle
% Trig Practice", not the full course name) — it shares the header line with
% the Name/Date blanks and a long title overlaps them. ──────────────────────
% \wsheader{Short Title} — student worksheet
\newcommand{\wsheader}[1]{%
  \pagestyle{fancy}\fancyhf{}%
  \fancyhead[L]{\textbf{\small #1}}%
  \fancyhead[R]{\small Name: \underline{\hspace{4.5cm}}~Date: \underline{\hspace{1.8cm}}}%
  \fancyfoot[C]{\thepage}%
  \renewcommand{\headrulewidth}{0.4pt}}

% \akheader{Short Topic} — answer key
\newcommand{\akheader}[1]{%
  \pagestyle{fancy}\fancyhf{}%
  \fancyhead[L]{\textbf{\small #1 --- Answer Key}}%
  \fancyhead[R]{\textit{\small For instructor/parent use}}%
  \fancyfoot[C]{\thepage}%
  \renewcommand{\headrulewidth}{0.4pt}}

% \ssheader{Topic} — skills summary / study guide
\newcommand{\ssheader}[1]{%
  \pagestyle{fancy}\fancyhf{}%
  \fancyhead[L]{\textbf{Skills Summary: #1}}%
  \fancyhead[R]{\small\textit{Study Guide \& Reference}}%
  \fancyfoot[C]{\thepage}%
  \renewcommand{\headrulewidth}{0.4pt}}

% ── Title blocks ────────────────────────────────────────────────────────────
% \wstitleblock{Title}{Course}{Date}
\newcommand{\wstitleblock}[3]{%
  \begin{center}
    \fittedtitle{#1}\\[0.3cm]
    {\large #2 \quad $\bullet$ \quad #3}
  \end{center}
  \noindent\rule{\linewidth}{0.4pt}\vspace{0.2cm}}

% \aktitleblock{Topic}{Course}{Date} — "Answer Key" is a subtitle BENEATH the
% topic, never appended to it: a long topic would wrap mid-phrase.
\newcommand{\aktitleblock}[3]{%
  \begin{center}
    \fittedtitle{#1}\\[0.15cm]
    {\Large\bfseries Answer Key}\\[0.25cm]
    {\large #2 \quad $\bullet$ \quad #3}
  \end{center}}

% \sstitleblock{Topic}
\newcommand{\sstitleblock}[1]{%
  \begin{center}
    {\LARGE\textbf{Skills Summary}}\\[0.2cm]
    {\large\textbf{#1}}\\[0.1cm]
    {\small\textit{Use this reference while completing your worksheet or when studying.}}
  \end{center}
  \vspace{0.1cm}
  \noindent\rule{\linewidth}{1.5pt}
  \vspace{0.3cm}}

% ── Study-guide color palette and boxes ─────────────────────────────────────
\definecolor{skillblue}{RGB}{30,100,180}
\definecolor{skillbluebg}{RGB}{235,244,255}
\definecolor{warnorange}{RGB}{200,90,0}
\definecolor{warnbg}{RGB}{255,243,230}
\definecolor{exgreen}{RGB}{20,120,60}
\definecolor{exgreenbg}{RGB}{230,248,238}

% Formula/rule box
\newmdenv[
  backgroundcolor=skillbluebg,
  linecolor=skillblue, linewidth=1.5pt,
  innertopmargin=6pt, innerbottommargin=6pt,
  innerleftmargin=10pt, innerrightmargin=10pt,
  skipabove=6pt, skipbelow=4pt
]{formulabox}

% Mini example box
\newmdenv[
  backgroundcolor=exgreenbg,
  linecolor=exgreen, linewidth=1pt,
  innertopmargin=5pt, innerbottommargin=5pt,
  innerleftmargin=10pt, innerrightmargin=10pt,
  skipabove=4pt, skipbelow=4pt
]{examplebox}

% Watch-out box
\newmdenv[
  backgroundcolor=warnbg,
  linecolor=warnorange, linewidth=1pt,
  innertopmargin=5pt, innerbottommargin=5pt,
  innerleftmargin=10pt, innerrightmargin=10pt,
  skipabove=4pt, skipbelow=4pt
]{watchoutbox}

% Skill section heading
\newcommand{\skillheading}[1]{%
  \vspace{0.4cm}
  {\large\textbf{\textcolor{skillblue}{#1}}}
  \vspace{0.1cm}
  \hrule height 1pt
  \vspace{0.2cm}
}

%   % figure-macros.tex — shipped geometric figure macros. SOURCE OF TRUTH: this
% file; references/latex-templates.md documents usage. Requires the tikz
% libraries loaded by worksheet-preamble.tex (calc, angles, quotes), so
% \input worksheet-preamble first.
%
% CONTRACT with the layout/prose checkers (do not break it):
%   - Macro names contain rt/tri/fig — that is how tests/check_layout.py and
%     tests/check_prose_consistency.py recognize a macro-style figure.
%   - Mandatory arguments are FLAT braces only (no nested {} inside an arg);
%     the checkers' arg scanners match [^{}]* per group.
%   - Value-bearing args carry exactly the printed label text; scale/styling
%     goes in the single optional [..] argument, which the checkers ignore
%     (so scale=1.2 is never mistaken for a figure value).
%
% Figure conventions (from SKILL.md / latex-templates.md): every number
% printed in a figure must come from the problem's verify JSON; vertices are
% A/B/C with sides a/b/c OPPOSITE them — the same convention verify.py's
% triangle type uses. Right triangles put the right angle at C, so c (= AB)
% is the hypotenuse — the SAME convention scripts/render_figures.py constructs
% (its RIGHT_ANGLE_VERTEX constant): \refrt sits beside renderer-built
% \probfig figures on the page, so the two sources must never disagree.
% tests/test_figure_convention.py enforces this by parsing the
% "% right angle mark" lines below — that trailing comment is load-bearing;
% keep it on every right-angle-mark \draw.
%
% ASCII only: pdflatex (the fallback engine) cannot typeset literal Unicode.

% ── \rtfig[scale]{bottom-leg label}{right-leg label}{hypotenuse label} ──────
% Right triangle, right angle at C (hypotenuse c = AB — the renderer's
% convention). A=(0,0), C=(4,0), B=(4,3): the 3-4-5 shape keeps every label
% clear of the drawn lines at default scale. Bottom leg = b (AC), right leg
% = a (CB), hypotenuse = c (AB) — name sides accordingly in the labels.
% Example: \rtfig{$8$}{$6$}{$x$}   or   \rtfig[0.9]{$b = 8$}{$a = 6$}{$c$}
\newcommand{\rtfig}[4][1.2]{%
  \begin{center}
  \begin{tikzpicture}[scale=#1]
    \coordinate (A) at (0,0);
    \coordinate (B) at (4,3);
    \coordinate (C) at (4,0);
    \draw[thick] (A) -- (C) -- (B) -- cycle;
    \draw (C) ++(-.25,0) -- ++(0,.25) -- ++(.25,0);  % right angle mark
    \node[below left] at (A) {$A$};
    \node[above right] at (B) {$B$};
    \node[below right] at (C) {$C$};
    \node[below] at (2,0) {#2};
    \node[right] at (4,1.5) {#3};
    \node[above left] at (2,1.5) {#4};
  \end{tikzpicture}
  \end{center}}

% ── \trifig[scale]{c}{b}{A-deg}{c-label}{b-label}{a-label}{angle-label} ─────
% General triangle, TO SCALE by construction (SAS): A at the origin, B on the
% x-axis at distance c, C computed from side b and angle A (TikZ trig takes
% degrees). Numeric args #2-#4 are the ACTUAL given values, so the drawing
% cannot disagree with the math; label args #5-#8 carry what is printed.
% Example: \trifig{7}{5}{34}{$c = 7$}{$b = 5$}{$a = ?$}{$34^\circ$}
\newcommand{\trifig}[8][0.6]{%
  \begin{center}
  \begin{tikzpicture}[scale=#1]
    \coordinate (A) at (0,0);
    \coordinate (B) at (#2,0);
    \coordinate (C) at ({#3*cos(#4)},{#3*sin(#4)});
    \draw[thick] (A) -- (B) -- (C) -- cycle;
    \node[below left]  at (A) {$A$};
    \node[below right] at (B) {$B$};
    \node[above]       at (C) {$C$};
    \node[below]       at ($(A)!0.5!(B)$) {#5};
    \node[above left]  at ($(A)!0.5!(C)$) {#6};
    \node[above right] at ($(B)!0.5!(C)$) {#7};
    \pic [draw, angle radius=7mm, angle eccentricity=1.6, "#8"] {angle = B--A--C};
  \end{tikzpicture}
  \end{center}}

% ── \refrt — the value-free reference figure ────────────────────────────────
% SKILL.md's figure-scope rule requires shared labelling conventions to live
% in ONE value-free reference figure placed with the directions; until now no
% template for it existed, so every generation invented one — and an invented
% figure that accidentally carried a number would itself trip the
% all-or-nothing figure rule it exists to satisfy. This macro has zero
% numerals in its output by construction (it takes no arguments and prints
% only letters), and the caption is baked in verbatim so the figure cannot be
% mistaken for a problem's givens. Right angle at C, so sides read: a = BC
% (opposite A), b = AC (opposite B), c = AB (opposite C, the hypotenuse) —
% matching render_figures.py's RIGHT_ANGLE_VERTEX, because this reference
% figure defines the convention students apply to every renderer-built
% figure on the sheet.
\newcommand{\refrt}{%
  \begin{center}
  \begin{tikzpicture}[scale=1.0]
    \coordinate (A) at (0,0);
    \coordinate (B) at (3.6,2.7);
    \coordinate (C) at (3.6,0);
    \draw[thick] (A) -- (C) -- (B) -- cycle;
    \draw (C) ++(-.25,0) -- ++(0,.25) -- ++(.25,0);  % right angle mark
    \node[below left] at (A) {$A$};
    \node[above right] at (B) {$B$};
    \node[below right] at (C) {$C$};
    \node[below] at (1.8,0) {$b$};
    \node[right] at (3.6,1.35) {$a$};
    \node[above left] at (1.8,1.35) {$c$};
    \pic [draw, angle radius=6mm, angle eccentricity=1.7, "$A$"] {angle = C--A--B};
  \end{tikzpicture}\\
  {\small\textit{How every triangle here is labelled. No values shown: use the
  numbers given in each problem.}}
  \end{center}}
          % when using the shipped figure macros
%
% geometry stays in the document, NOT here: worksheets use margin=1in
% (top/bottom 0.75in) while study guides use margin=0.85in (top/bottom 0.7in),
% and geometry options cannot be re-issued once the package is loaded.
%
% ASCII only in this file: pdflatex (the fallback engine) cannot typeset
% literal Unicode symbols, so documents must compile identically under both
% tectonic and pdflatex.

% ── Required packages (superset of the worksheet and study-guide shells; a
% worksheet loading mdframed it does not use costs nothing, one shared file
% that cannot drift is worth it) ─────────────────────────────────────────────
\usepackage{amsmath, amssymb}
\usepackage{tikz, pgfplots}
\usepackage{enumitem, fancyhdr, multicol, array, booktabs}
\usepackage{mdframed, xcolor}
\pgfplotsset{compat=1.18}
\usetikzlibrary{calc, angles, quotes}   % needed by the geometric figure macros

% ── Problem numbering ───────────────────────────────────────────────────────
\newcounter{prob}
\newcommand{\problem}[1]{\stepcounter{prob}\vspace{0.4cm}\noindent\textbf{\large\theprob.}\quad #1\vspace{0.3cm}}

% ── Shrink-to-fit title: \LARGE when it fits on one line, otherwise scale
% down to the text width — never enlarges, never wraps mid-phrase. (\resizebox
% comes from graphicx, which tikz already loads.) ────────────────────────────
\newlength{\fittedw}
\newcommand{\fittedtitle}[1]{%
  \settowidth{\fittedw}{\LARGE\bfseries #1}%
  \ifdim\fittedw>\linewidth\resizebox{\linewidth}{!}{\bfseries #1}%
  \else{\LARGE\bfseries #1}\fi}

% ── Running headers. Keep the left title SHORT (<= ~28 chars, e.g. "Triangle
% Trig Practice", not the full course name) — it shares the header line with
% the Name/Date blanks and a long title overlaps them. ──────────────────────
% \wsheader{Short Title} — student worksheet
\newcommand{\wsheader}[1]{%
  \pagestyle{fancy}\fancyhf{}%
  \fancyhead[L]{\textbf{\small #1}}%
  \fancyhead[R]{\small Name: \underline{\hspace{4.5cm}}~Date: \underline{\hspace{1.8cm}}}%
  \fancyfoot[C]{\thepage}%
  \renewcommand{\headrulewidth}{0.4pt}}

% \akheader{Short Topic} — answer key
\newcommand{\akheader}[1]{%
  \pagestyle{fancy}\fancyhf{}%
  \fancyhead[L]{\textbf{\small #1 --- Answer Key}}%
  \fancyhead[R]{\textit{\small For instructor/parent use}}%
  \fancyfoot[C]{\thepage}%
  \renewcommand{\headrulewidth}{0.4pt}}

% \ssheader{Topic} — skills summary / study guide
\newcommand{\ssheader}[1]{%
  \pagestyle{fancy}\fancyhf{}%
  \fancyhead[L]{\textbf{Skills Summary: #1}}%
  \fancyhead[R]{\small\textit{Study Guide \& Reference}}%
  \fancyfoot[C]{\thepage}%
  \renewcommand{\headrulewidth}{0.4pt}}

% ── Title blocks ────────────────────────────────────────────────────────────
% \wstitleblock{Title}{Course}{Date}
\newcommand{\wstitleblock}[3]{%
  \begin{center}
    \fittedtitle{#1}\\[0.3cm]
    {\large #2 \quad $\bullet$ \quad #3}
  \end{center}
  \noindent\rule{\linewidth}{0.4pt}\vspace{0.2cm}}

% \aktitleblock{Topic}{Course}{Date} — "Answer Key" is a subtitle BENEATH the
% topic, never appended to it: a long topic would wrap mid-phrase.
\newcommand{\aktitleblock}[3]{%
  \begin{center}
    \fittedtitle{#1}\\[0.15cm]
    {\Large\bfseries Answer Key}\\[0.25cm]
    {\large #2 \quad $\bullet$ \quad #3}
  \end{center}}

% \sstitleblock{Topic}
\newcommand{\sstitleblock}[1]{%
  \begin{center}
    {\LARGE\textbf{Skills Summary}}\\[0.2cm]
    {\large\textbf{#1}}\\[0.1cm]
    {\small\textit{Use this reference while completing your worksheet or when studying.}}
  \end{center}
  \vspace{0.1cm}
  \noindent\rule{\linewidth}{1.5pt}
  \vspace{0.3cm}}

% ── Study-guide color palette and boxes ─────────────────────────────────────
\definecolor{skillblue}{RGB}{30,100,180}
\definecolor{skillbluebg}{RGB}{235,244,255}
\definecolor{warnorange}{RGB}{200,90,0}
\definecolor{warnbg}{RGB}{255,243,230}
\definecolor{exgreen}{RGB}{20,120,60}
\definecolor{exgreenbg}{RGB}{230,248,238}

% Formula/rule box
\newmdenv[
  backgroundcolor=skillbluebg,
  linecolor=skillblue, linewidth=1.5pt,
  innertopmargin=6pt, innerbottommargin=6pt,
  innerleftmargin=10pt, innerrightmargin=10pt,
  skipabove=6pt, skipbelow=4pt
]{formulabox}

% Mini example box
\newmdenv[
  backgroundcolor=exgreenbg,
  linecolor=exgreen, linewidth=1pt,
  innertopmargin=5pt, innerbottommargin=5pt,
  innerleftmargin=10pt, innerrightmargin=10pt,
  skipabove=4pt, skipbelow=4pt
]{examplebox}

% Watch-out box
\newmdenv[
  backgroundcolor=warnbg,
  linecolor=warnorange, linewidth=1pt,
  innertopmargin=5pt, innerbottommargin=5pt,
  innerleftmargin=10pt, innerrightmargin=10pt,
  skipabove=4pt, skipbelow=4pt
]{watchoutbox}

% Skill section heading
\newcommand{\skillheading}[1]{%
  \vspace{0.4cm}
  {\large\textbf{\textcolor{skillblue}{#1}}}
  \vspace{0.1cm}
  \hrule height 1pt
  \vspace{0.2cm}
}

\wsheader{The Coordinate Grid}

\begin{document}
\wstitleblock{Reading and Plotting Points on a Grid}{Grade 4--5}{}

\begin{center}\small\textit{Reading ordered pairs off a first-quadrant grid, plotting new
points on it, and turning two number patterns into points}\end{center}

\noindent\small\textbf{How to use this sheet.}
\begin{itemize}[label={$\bullet$}, leftmargin=1.1cm, itemsep=1pt, topsep=2pt]
  \item An \emph{ordered pair} is written $(x, y)$. The first number tells you how far to
        travel \textbf{across} from the corner marked $0$, and the second tells you how far
        to travel \textbf{up}. Across first, then up --- every time, in that order.
  \item Every grid on this sheet starts at $0$ in the bottom-left corner and counts by ones
        in both directions.
  \item Say the journey out loud as you go: ``start at zero, across four, up eight.'' The
        one mistake this sheet keeps coming back to is doing those two steps in the wrong
        order.
  \item Use a ruler or the edge of a card to trace across and up. It is much easier than
        counting squares by eye.
\end{itemize}
\normalsize

\bigskip

\problem[0.4cm]{\textbf{Read point $A$.} Start at the corner. Trace across to $A$, then up to $A$.

\par\smallskip
\noindent\begin{minipage}[t]{0.5\linewidth}\vspace{0pt}
\begin{tikzpicture}[scale=0.5]
  \draw[gray!35, very thin] (0,0) grid (10,10);
  \draw[->, thick] (0,0) -- (10.7,0) node[right, font=\small] {$x$};
  \draw[->, thick] (0,0) -- (0,10.7) node[above, font=\small] {$y$};
  \foreach \t in {1,...,10} {\node[below, font=\tiny] at (\t,0) {\t}; \node[left, font=\tiny] at (0,\t) {\t};}
  \node[below left, font=\tiny] at (0,0) {0};
  \fill[blue!80] (2,5) circle (3pt);
  \node[above right, font=\small] at (2,5) {$A$};
  \fill[blue!80] (6,3) circle (3pt);
  \node[above right, font=\small] at (6,3) {$B$};
  \fill[blue!80] (4,8) circle (3pt);
  \node[above right, font=\small] at (4,8) {$C$};
  \fill[blue!80] (9,1) circle (3pt);
  \node[above right, font=\small] at (9,1) {$D$};
  \fill[blue!80] (0,6) circle (3pt);
  \node[above right, font=\small] at (0,6) {$E$};
  \fill[blue!80] (7,7) circle (3pt);
  \node[above right, font=\small] at (7,7) {$F$};
\end{tikzpicture}
\end{minipage}\hfill
\begin{minipage}[t]{0.46\linewidth}\vspace{0pt}
Point $A$ is marked in blue with five other points.
\par\smallskip
How far across is $A$ from $0$?~\ansblank
\par\smallskip
How far up is $A$ from $0$?~\ansblank
\par\smallskip
Write it as an ordered pair: $A = ($~\ansblank$,$~\ansblank$)$
\par\vspace*{2.4cm}
\end{minipage}}

\problem[0.4cm]{\textbf{Read point $B$.} Same grid, same method.

\par\smallskip
\noindent\begin{minipage}[t]{0.5\linewidth}\vspace{0pt}
\begin{tikzpicture}[scale=0.5]
  \draw[gray!35, very thin] (0,0) grid (10,10);
  \draw[->, thick] (0,0) -- (10.7,0) node[right, font=\small] {$x$};
  \draw[->, thick] (0,0) -- (0,10.7) node[above, font=\small] {$y$};
  \foreach \t in {1,...,10} {\node[below, font=\tiny] at (\t,0) {\t}; \node[left, font=\tiny] at (0,\t) {\t};}
  \node[below left, font=\tiny] at (0,0) {0};
  \fill[blue!80] (2,5) circle (3pt);
  \node[above right, font=\small] at (2,5) {$A$};
  \fill[blue!80] (6,3) circle (3pt);
  \node[above right, font=\small] at (6,3) {$B$};
  \fill[blue!80] (4,8) circle (3pt);
  \node[above right, font=\small] at (4,8) {$C$};
  \fill[blue!80] (9,1) circle (3pt);
  \node[above right, font=\small] at (9,1) {$D$};
  \fill[blue!80] (0,6) circle (3pt);
  \node[above right, font=\small] at (0,6) {$E$};
  \fill[blue!80] (7,7) circle (3pt);
  \node[above right, font=\small] at (7,7) {$F$};
\end{tikzpicture}
\end{minipage}\hfill
\begin{minipage}[t]{0.46\linewidth}\vspace{0pt}
Across~\ansblank \qquad Up~\ansblank
\par\smallskip
$B = ($~\ansblank$,$~\ansblank$)$
\par\smallskip
\small Now write $C$ and $D$ as ordered pairs too, for practice:
\par\smallskip
$C = ($ \underline{\hspace{0.8cm}} $,$ \underline{\hspace{0.8cm}} $)$ \qquad
$D = ($ \underline{\hspace{0.8cm}} $,$ \underline{\hspace{0.8cm}} $)$\normalsize
\par\vspace*{2.4cm}
\end{minipage}}

\problem[0.4cm]{\textbf{A point that sits on a line.} Point $E$ is not out in the middle of the grid --- it is sitting right on one of the two arrows.

\par\smallskip
\noindent\begin{minipage}[t]{0.5\linewidth}\vspace{0pt}
\begin{tikzpicture}[scale=0.5]
  \draw[gray!35, very thin] (0,0) grid (10,10);
  \draw[->, thick] (0,0) -- (10.7,0) node[right, font=\small] {$x$};
  \draw[->, thick] (0,0) -- (0,10.7) node[above, font=\small] {$y$};
  \foreach \t in {1,...,10} {\node[below, font=\tiny] at (\t,0) {\t}; \node[left, font=\tiny] at (0,\t) {\t};}
  \node[below left, font=\tiny] at (0,0) {0};
  \fill[blue!80] (2,5) circle (3pt);
  \node[above right, font=\small] at (2,5) {$A$};
  \fill[blue!80] (6,3) circle (3pt);
  \node[above right, font=\small] at (6,3) {$B$};
  \fill[blue!80] (4,8) circle (3pt);
  \node[above right, font=\small] at (4,8) {$C$};
  \fill[blue!80] (9,1) circle (3pt);
  \node[above right, font=\small] at (9,1) {$D$};
  \fill[blue!80] (0,6) circle (3pt);
  \node[above right, font=\small] at (0,6) {$E$};
  \fill[blue!80] (7,7) circle (3pt);
  \node[above right, font=\small] at (7,7) {$F$};
\end{tikzpicture}
\end{minipage}\hfill
\begin{minipage}[t]{0.46\linewidth}\vspace{0pt}
How far across is $E$ from the corner?~\ansblank
\par\smallskip
How far up is $E$?~\ansblank
\par\smallskip
$E = ($~\ansblank$,$~\ansblank$)$
\par\smallskip
\small Which arrow is $E$ sitting on? \underline{\hspace{3cm}}\normalsize
\par\vspace*{2.2cm}
\end{minipage}}

\problem[0.4cm]{\textbf{Plot two points, then measure across.} Plot $P(3, 6)$ and $Q(8, 6)$ on the empty grid, label them, and join them with a straight line.

\par\smallskip
\noindent\begin{minipage}[t]{0.5\linewidth}\vspace{0pt}
\begin{tikzpicture}[scale=0.5]
  \draw[gray!35, very thin] (0,0) grid (10,10);
  \draw[->, thick] (0,0) -- (10.7,0) node[right, font=\small] {$x$};
  \draw[->, thick] (0,0) -- (0,10.7) node[above, font=\small] {$y$};
  \foreach \t in {1,...,10} {\node[below, font=\tiny] at (\t,0) {\t}; \node[left, font=\tiny] at (0,\t) {\t};}
  \node[below left, font=\tiny] at (0,0) {0};
\end{tikzpicture}
\end{minipage}\hfill
\begin{minipage}[t]{0.46\linewidth}\vspace{0pt}
Both points are the same height up. So the line you drew is flat.
\par\smallskip
$x$ of $Q$~\ansblank \quad $-$ \quad $x$ of $P$~\ansblank
\par\smallskip
How many squares long is $\overline{PQ}$?~\ansblank
\par\smallskip
\small Why did you not need the $y$-coordinates to answer that?
\underline{\hspace{4cm}}\normalsize
\par\vspace*{2.4cm}
\end{minipage}}

\problem[0.4cm]{\textbf{Now measure up.} Plot $R(4, 2)$ and $S(4, 9)$, label them, and join them.

\par\smallskip
\noindent\begin{minipage}[t]{0.5\linewidth}\vspace{0pt}
\begin{tikzpicture}[scale=0.5]
  \draw[gray!35, very thin] (0,0) grid (10,10);
  \draw[->, thick] (0,0) -- (10.7,0) node[right, font=\small] {$x$};
  \draw[->, thick] (0,0) -- (0,10.7) node[above, font=\small] {$y$};
  \foreach \t in {1,...,10} {\node[below, font=\tiny] at (\t,0) {\t}; \node[left, font=\tiny] at (0,\t) {\t};}
  \node[below left, font=\tiny] at (0,0) {0};
\end{tikzpicture}
\end{minipage}\hfill
\begin{minipage}[t]{0.46\linewidth}\vspace{0pt}
This time both points are the same distance across, so the line is upright.
\par\smallskip
$y$ of $S$~\ansblank \quad $-$ \quad $y$ of $R$~\ansblank
\par\smallskip
How many squares long is $\overline{RS}$?~\ansblank
\par\smallskip
\small Which coordinate is the same for both points?
\underline{\hspace{3cm}}\normalsize
\par\vspace*{2.4cm}
\end{minipage}}

\problem[0.4cm]{\textbf{How much higher?} Look again at the first grid. Point $C$ is at $(4, 8)$ and point $B$ is at $(6, 3)$.

\par\smallskip
\noindent\begin{minipage}[t]{0.5\linewidth}\vspace{0pt}
\begin{tikzpicture}[scale=0.5]
  \draw[gray!35, very thin] (0,0) grid (10,10);
  \draw[->, thick] (0,0) -- (10.7,0) node[right, font=\small] {$x$};
  \draw[->, thick] (0,0) -- (0,10.7) node[above, font=\small] {$y$};
  \foreach \t in {1,...,10} {\node[below, font=\tiny] at (\t,0) {\t}; \node[left, font=\tiny] at (0,\t) {\t};}
  \node[below left, font=\tiny] at (0,0) {0};
  \fill[blue!80] (2,5) circle (3pt);
  \node[above right, font=\small] at (2,5) {$A$};
  \fill[blue!80] (6,3) circle (3pt);
  \node[above right, font=\small] at (6,3) {$B$};
  \fill[blue!80] (4,8) circle (3pt);
  \node[above right, font=\small] at (4,8) {$C$};
  \fill[blue!80] (9,1) circle (3pt);
  \node[above right, font=\small] at (9,1) {$D$};
  \fill[blue!80] (0,6) circle (3pt);
  \node[above right, font=\small] at (0,6) {$E$};
  \fill[blue!80] (7,7) circle (3pt);
  \node[above right, font=\small] at (7,7) {$F$};
\end{tikzpicture}
\end{minipage}\hfill
\begin{minipage}[t]{0.46\linewidth}\vspace{0pt}
``How much higher'' is a question about going \emph{up}, so it is the second number
in each pair you need.
\par\smallskip
$y$ of $C$~\ansblank \quad $-$ \quad $y$ of $B$~\ansblank
\par\smallskip
$C$ is~\ansblank squares higher than $B$.
\par\smallskip
\small Careful: $B$ is further across than $C$. Does that change how much higher $C$ is?
\underline{\hspace{2.5cm}}\normalsize
\par\vspace*{2.2cm}
\end{minipage}}

\problem[0.4cm]{\textbf{Two rules make points.} Rule A: \emph{start at $0$ and add $1$ each time.}
Rule B: \emph{start at $0$ and add $2$ each time.} Take the first number from Rule A and
the second from Rule B, and you get an ordered pair.

\par\smallskip
\begin{center}
\renewcommand{\arraystretch}{1.3}
\begin{tabular}{|c|c|c|c|c|c|}
\hline
Rule A ($x$) & $0$ & $1$ & $2$ & $3$ & $4$ \\ \hline
Rule B ($y$) & $0$ & $2$ & \ansblank & \ansblank & \ansblank \\ \hline
\end{tabular}
\end{center}

\par\smallskip
\noindent\begin{minipage}[t]{0.5\linewidth}\vspace{0pt}
\begin{tikzpicture}[scale=0.5]
  \draw[gray!35, very thin] (0,0) grid (10,10);
  \draw[->, thick] (0,0) -- (10.7,0) node[right, font=\small] {$x$};
  \draw[->, thick] (0,0) -- (0,10.7) node[above, font=\small] {$y$};
  \foreach \t in {1,...,10} {\node[below, font=\tiny] at (\t,0) {\t}; \node[left, font=\tiny] at (0,\t) {\t};}
  \node[below left, font=\tiny] at (0,0) {0};
\end{tikzpicture}
\end{minipage}\hfill
\begin{minipage}[t]{0.46\linewidth}\vspace{0pt}
Plot all five ordered pairs on the grid and join them in order.
\par\smallskip
$y$ when $x = 3$:~\ansblank
\par\smallskip
$y$ when $x = 4$:~\ansblank
\par\smallskip
\small Finish the sentence: every $y$ is always \underline{\hspace{1.5cm}} times the $x$
beside it.\normalsize
\par\vspace*{2.2cm}
\end{minipage}}

\problem[0.4cm]{\textbf{Keep the pattern going.} Use the same two rules as problem 7. What is the next ordered pair after $(4, 8)$?

\par\smallskip
\noindent\begin{minipage}[t]{0.5\linewidth}\vspace{0pt}
\begin{tikzpicture}[scale=0.5]
  \draw[gray!35, very thin] (0,0) grid (10,10);
  \draw[->, thick] (0,0) -- (10.7,0) node[right, font=\small] {$x$};
  \draw[->, thick] (0,0) -- (0,10.7) node[above, font=\small] {$y$};
  \foreach \t in {1,...,10} {\node[below, font=\tiny] at (\t,0) {\t}; \node[left, font=\tiny] at (0,\t) {\t};}
  \node[below left, font=\tiny] at (0,0) {0};
\end{tikzpicture}
\end{minipage}\hfill
\begin{minipage}[t]{0.46\linewidth}\vspace{0pt}
Rule A adds $1$, so the next $x$ is~\ansblank
\par\smallskip
Rule B adds $2$, so the next $y$ is~\ansblank
\par\smallskip
The next point is $($~\ansblank$,$~\ansblank$)$. Plot it on the grid.
\par\smallskip
\small Does the new point sit on the same straight line as the others?
\underline{\hspace{2.5cm}}\normalsize
\par\vspace*{2.4cm}
\end{minipage}}

\problem[0.4cm]{\textbf{A graph that tells a story.} Each point shows how many books Maya read in
one week. Week $1$ is at $(1, 3)$, and the rest follow.

\par\smallskip
\noindent\begin{minipage}[t]{0.5\linewidth}\vspace{0pt}
\begin{tikzpicture}[scale=0.5]
  \draw[gray!35, very thin] (0,0) grid (10,10);
  \draw[->, thick] (0,0) -- (10.7,0) node[right, font=\small] {$x$};
  \draw[->, thick] (0,0) -- (0,10.7) node[above, font=\small] {$y$};
  \foreach \t in {1,...,10} {\node[below, font=\tiny] at (\t,0) {\t}; \node[left, font=\tiny] at (0,\t) {\t};}
  \node[below left, font=\tiny] at (0,0) {0};
  \fill[blue!80] (1,3) circle (3pt);
  \fill[blue!80] (2,5) circle (3pt);
  \fill[blue!80] (3,4) circle (3pt);
  \fill[blue!80] (4,8) circle (3pt);
  \fill[blue!80] (5,5) circle (3pt);
\end{tikzpicture}
\end{minipage}\hfill
\begin{minipage}[t]{0.46\linewidth}\vspace{0pt}
Read the height of each point and write it down.
\par\smallskip
Wk 1~\ansblank \ Wk 2~\ansblank \ Wk 3~\ansblank \ Wk 4~\ansblank \ Wk 5~\ansblank
\par\smallskip
Books altogether:~\ansblank
\par\smallskip
Share them equally over the five weeks: mean~\ansblank books per week
\par\smallskip
\small Which week was Maya's best? \underline{\hspace{2cm}}\normalsize
\par\vspace*{2cm}
\end{minipage}}

\problem[0.4cm]{\textbf{Explain it to a friend.} Nia says: ``$(3, 7)$ and $(7, 3)$ must be the same
point, because they use the same two numbers.'' Both points are marked on the grid.

\par\smallskip
\noindent\begin{minipage}[t]{0.5\linewidth}\vspace{0pt}
\begin{tikzpicture}[scale=0.5]
  \draw[gray!35, very thin] (0,0) grid (10,10);
  \draw[->, thick] (0,0) -- (10.7,0) node[right, font=\small] {$x$};
  \draw[->, thick] (0,0) -- (0,10.7) node[above, font=\small] {$y$};
  \foreach \t in {1,...,10} {\node[below, font=\tiny] at (\t,0) {\t}; \node[left, font=\tiny] at (0,\t) {\t};}
  \node[below left, font=\tiny] at (0,0) {0};
  \fill[blue!80] (3,7) circle (3pt);
  \node[above left, font=\small] at (3,7) {$(3,7)$};
  \fill[blue!80] (7,3) circle (3pt);
  \node[above right, font=\small] at (7,3) {$(7,3)$};
\end{tikzpicture}
\end{minipage}\hfill
\begin{minipage}[t]{0.46\linewidth}\vspace{0pt}
Trace the journey to each one with your finger, saying it out loud.
\par\smallskip
Then write two or three sentences telling Nia why the two points are not the same, and what
the \emph{order} of the two numbers is for.

\par\smallskip
\noindent\rule{\linewidth}{0.4pt}

\vspace{0.9cm}
\noindent\rule{\linewidth}{0.4pt}

\vspace{0.9cm}
\noindent\rule{\linewidth}{0.4pt}

\vspace{0.9cm}
\noindent\rule{\linewidth}{0.4pt}
\par\vspace*{0.2cm}
\noansline
\end{minipage}}

\end{document}
